\chapter{Adaptación de investigación a Philips Azurion}
\section{Identidad del equipo y alcance del perfil}
El trabajo reciente utiliza un RDSR que declara fabricante Philips y modelo Azurion. El perfil incorpora controles de identidad y políticas específicas de lectura y configuración. Esto permite trabajar con las particularidades observadas sin cambiar el informe original ni extender una suposición a cualquier equipo de la misma marca.

La denominación Azurion identifica una familia, no un único espectro. La instalación puede diferir por tubo, configuración de limitador, versión, mantenimiento y programa clínico. Para construir un perfil físico se necesita identificar esos elementos. Una configuración denominada ``Azurion'' que conserva factores sintéticos sigue siendo de investigación.

\section{Qué aporta la documentación del fabricante}
La documentación oficial consultada para Release 1.2 describe una filtración mínima del conjunto tubo--carcasa equivalente a \num{2.5} mm Al. También identifica filtros añadidos de \num{0.1}, \num{0.4} y \num{0.9} mm Cu, cada uno con 1 mm Al, para el limitador de haz especificado \citep{philipsifu}. La nota 41 del documento ClarityIQ confirma la presencia del aluminio acompañante al cobre \citep{clarity}.

Esta información fundamenta conservar ambos materiales. No determina la filtración fija efectiva total de la instalación ni justifica por sí sola \num{3.1} mm Al. Los espesores equivalentes se refieren a condiciones de calidad del haz declaradas; no representan una equivalencia espectral universal. La versión y el limitador descritos deben contrastarse con los de la sala.

Por tanto, el parámetro de \num{3.1} mm Al de la configuración se clasifica como hipótesis. El ángulo de $12^\circ$ pertenece exclusivamente al dataset histórico. Para el MRC 200 0407 identificado se documentan $11^\circ$; véase el capítulo de caracterización actualizado. El blanco de tungsteno es parte del modelo espectral utilizado y requiere esa misma coherencia de identificación.

\section{Inventario del RDSR de trabajo}
La inspección agregada repetida para esta edición obtiene 101 eventos y tres grupos de filtración añadida. La tabla~\ref{tab:groups} no contiene datos identificativos del procedimiento ni pretende representar la frecuencia de uso del equipo en general.
\begin{table}[htbp]\centering
\caption{Filtración añadida observada en el RDSR de trabajo.}\label{tab:groups}
\begin{tabular}{lrrr}\toprule
Estado & Al añadido (mm) & Cu añadido (mm) & Eventos\\\midrule
Filtro no insertado, explícito & 0 & 0 & 30\\
Pila Al + Cu & 1 & \num{0.1} & 47\\
Pila Al + Cu & 1 & \num{0.4} & 24\\\midrule
Total & & & 101\\\bottomrule
\end{tabular}
\fuente{Inspección de \codigo{0_anon.dcm} mediante \codigo{inspect_azurion}; salida agregada conservada, sin exportar el DICOM.}
\end{table}

El error original ``material filtrante indeterminado'' no demostraba que el tubo careciera de aluminio. Se producía al tratar una declaración de filtro no insertado como si siempre debiera contener material y espesor. La corrección de lectura resuelve esa condición, manteniendo desconocido un filtro realmente no informado.

Hay un evento con kerma no nulo y las tres distancias SID, SOD y SRP nulas. Una distancia cero no permite aplicar el modelo geométrico. No se puede deducir la geometría de ese evento únicamente porque los anteriores tengan valores válidos o porque pertenezca a la misma sala. La sustitución por un valor típico necesita declararse como hipótesis.

\section{Mapeo de aluminio fijo y añadido}
La política \codigo{fixed_plus_event} selecciona una familia según
\begin{equation}\label{eq:al}
a_{\mathrm{total},e}=a_{\mathrm{fijo}}+a_{\text{añadido},e},\qquad
c_e=c_{\text{añadido},e}.
\end{equation}
Con la hipótesis $a_{\mathrm{fijo}}=\num{3.1}$ mm, los eventos explícitamente sin filtro añadido utilizan \num{3.1} mm Al y 0 Cu; los demás utilizan \num{4.1} mm Al y conservan su cobre. El audit registra los tres valores de aluminio: fijo, añadido y familia elegida.
\begin{figure}[htbp]\centering\begin{tikzpicture}[font=\small,x=1cm,y=1cm]
\node[punto,label=above:{Foco W}]at(.5,1.2){};
\fill[azul!8](.5,1.2)--(13.5,2.1)--(13.5,.3)--cycle;
\draw[rayo](.5,1.2)--(13.5,1.2);
\fill[azul!35](3,.1)rectangle(3.4,2.3);
\fill[gray!35](6.4,.1)rectangle(6.8,2.3);
\fill[ocre!55](9.2,.1)rectangle(9.6,2.3);
\node[rotulo,align=center]at(3.2,3){Al fijo efectivo\\$a_{\rm fijo}$};
\node[rotulo,align=center]at(6.6,3){Al añadido\\$a_e$};
\node[rotulo,align=center]at(9.4,3){Cu añadido\\$c_e$};
\node[nota,align=center]at(3.2,-.5){Hipótesis del perfil\\histórico: $3{,}1$ mm};
\node[nota,align=center]at(6.6,-.5){Desde el evento\\$0$ o $1$ mm};
\node[nota,align=center]at(9.4,-.5){Desde el evento\\$0$; $0{,}1$; $0{,}4$ mm};
\node[proceso,text width=12.6cm,anchor=north west]at(.2,-1.4){\textbf{Espectro del evento:} $Q=(\mathrm{kVp},\mathrm W,\theta,a_{\rm fijo}+a_e,c_e)$.\\``No filter'' significa ausencia declarada de filtro añadido; no elimina los componentes fijos del tubo.};
\end{tikzpicture}

\caption{Construcción de la filtración del espectro de investigación. Los espesores de Al no se convierten a Cu equivalente.}\end{figure}

La clave heredada \codigo{inherent_al_mm} representa bajo esta política la hipótesis de Al fijo efectivo, aunque su nombre sugiera una interpretación más restringida. Es necesario documentarlo para no sumar otra vez componentes ya incluidos. Al generar una fila del dataset, el argumento de Al es el total de esa fila; no se añade de nuevo el Al del evento.

Solo se seleccionan familias presentes, con tolerancia numérica para igualdad de flotantes. No se interpola entre familias ni se modifica una tabla original para alojar una combinación nueva. El proveedor comprueba también kVp, Cu y lado de campo local.

\section{Por qué la tabla clínica era insuficiente}
Los 24 eventos de \num{0.4} mm Cu seleccionan, con la hipótesis anterior, la familia de \num{4.1} mm Al. La tabla clínica original de esa familia termina en \num{0.2} mm Cu. El fallo indica falta de cobertura, no un error aritmético ni una imposibilidad de que ese filtro exista en el equipo.

El RDSR contiene además un evento de 46 kVp, por debajo del límite clínico publicado de 50 kVp. Aunque el eje de filtración se resolviera, ese kVp seguiría fuera de la tabla clínica. El campo local introduce una tercera condición independiente: su lado equivalente debe quedar entre 5 y 35 cm.

Usar constantes históricas en los eventos restantes permitiría obtener un número, pero mezclaría dos modelos físicos y ocultaría un cambio de calidad. La decisión implementada para Azurion es mantener un proveedor espectral derivado con procedencia propia y bloquear lo que aún quede fuera, sin sustitución automática por constantes.

\section{Perfil de investigación y plantilla de caracterización}
La configuración de investigación activa \codigo{allow_research_assumptions}. Utiliza distancias configuradas de 810 mm foco--isocentro, \num{1044.3} mm foco--detector y 660 mm foco--referencia cuando se autoriza reemplazar una distancia ausente o no positiva. Son valores de trabajo, derivados de valores habituales válidos del archivo examinado; no son medidas del evento inválido ni especificaciones universales Azurion.

El perfil conserva las distancias válidas de cada evento. También puede tomar las coordenadas de mesa del primer evento como referencia relativa. Esto permite estudiar cambios respecto al inicio, pero no determina dónde se sitúan cabeza, tronco o pelvis respecto al isocentro. El registro anatómico continúa pendiente.

La plantilla de caracterización no permite esas hipótesis por defecto y exige estados de filtración, geometría y calibración caracterizadas. Su función es hacer visible la preparación necesaria. El cambio de indicadores debe acompañarse de nuevos parámetros y evidencia; una modificación del JSON por sí sola no transforma el modelo de investigación en uno validado.

\section{Alcance real de la compatibilidad conseguida}
Se ha conseguido leer las tres categorías de filtros observadas, seleccionar Al total de forma trazable y disponer de una tabla calculada que incluye sus combinaciones de calidad bajo hipótesis. El perfil detecta fabricante/modelo incompatible y no ignora silenciosamente los errores de dominio.

No se ha demostrado que cualquier RDSR Azurion pueda procesarse sin incidencias. Podrían aparecer materiales no soportados, cuñas, campos fuera de rango, otro limitador, barridos o datos geométricos insuficientes. La tabla histórica de 390 filas no contiene \num{0.9} mm Cu, aunque ese filtro figure entre las opciones documentadas del fabricante: ese dataset se construyó para las calidades de trabajo. La alternativa posterior de 910 filas amplía este eje a 0,9 mm Cu, sin demostrar compatibilidad con toda la familia Azurion.

La respuesta científicamente correcta a la pregunta de si ``ya trabaja con Azurion'' es condicional: hay soporte específico de investigación, con lectura y cálculo ampliados, pero faltan caracterización de instalación, registro y validación independiente. La condición debe acompañar al resultado y no quedar solo en una nota de desarrollo.
\FloatBarrier
